\SourceLeaf{242}%
\Needspace{6\baselineskip}\begin{center}\textsc{Gauss} an \textsc{Gerling}. Göttingen, 4. Februar 1844.\end{center}

\mbox{.\ .\ .\ .\ .\ .\ .\ .} So kam es, dass ich vergessen habe, einen \textls{anderen} Punkt Ihres Briefs zu beantworten, an den mich erst heute ein ganz zufälliger Anlass wieder erinnert hat.

Der Titel des Buchs meines alten Universitätsfreundes \textsc{Varkas} (i. e. Wolfgang) von \textsc{Bolyai}, des Vaters, lautet wörtlich wie folgt.

Tentamen juventutem studiosam in elementa matheseos purae, elementaris ac sublimioris, methodo intuitiva, evidenti\"aque huic propria, introducendi cum Appendice triplici. Auctore Professore Matheseos et Physices Chemiaeque Publ. Ordinario. Tomus primus. Mar[o]s Vásárhelyini 1832. Typis Collegii Reformatorum per Josephum, et Simeonem Kali de felsö Vist. LXVI u. 502 Seiten 8\textsuperscript{vo}. Tomus secundus 1833. XVI u. 400 S., viele Kupfertafeln. Den Namen des Verfassers finde ich nicht angegeben.

Der eine Anhang im ersten Band trägt den Titel

Appendix scientiam spatii absolute veram exhibens: a veritate aut falsitate Axiomatis XI Euclidei (a priori haud unquam decidenda) independentem; adjecta ad casum falsitatis, quadratura circuli geometrica. Auctore \textsc{Johanne} \textsc{Bolyai} de eadem\HistoricalNote{*)}{*)}{\emph{of that ilk}, wie es gewöhnlich bei \textsc{Walter Scott} heißt. — [\emph{Vgl. auch S.} 225.]}, Geometrarum in exercitu Caesareo Regio Austriaco Castrensium Capitaneo.

Dieser Anhang hat eine eigene Seitenzählung (S. 1---26). Da auf dem Titelblatt weder Ort noch Jahr angegeben sind, vermute ich, dass diese Schrift, \textls{um die es hier gerade geht}, nur als Anhang zum Werk des Vaters veröffentlicht wurde.

Übrigens hat in den letzten Jahrzehnten ein Russe (\textsc{Lobatschewski}, Staatsrat und Professor in Kasan) einen ähnlichen Weg eingeschlagen. Er nennt die nichteuklidische Geometrie imaginäre Geometrie, so wie Ihr ehemaliger Kollege [\textsc{Schweikart}] sie Astralgeometrie nannte. Er hat darüber auf \textls{Russisch} viele sehr umfangreiche Abhandlungen veröffentlicht, meist in den \gausscyr{Записки Казанскаго Университета}, den Memoiren der Kasaner Universität, zum Teil auch in eigenen Broschüren. Ich glaube, dass ich sie alle besitze, habe ihre gründliche Lektüre aber verschoben,
